% =========================================================================
% ABSTRACT
% =========================================================================

\chapter*{Abstract}
\addcontentsline{toc}{chapter}{Abstract}

Many of the systems that characterize the world around us are either invisible, or extremely complex because dominated by countless variables that are often dependent on one another, or difficult to predict, or both these conditions together. They often move beneath the surface, unfold over timescales longer than a human life, and cannot be understood from any single moment. The difficulty is not a lack of data but the gap between data and understanding, a gap that conventional static visualisations often widen for non-specialists rather than close.

This thesis asks to what extent machine-learning-driven, experiential visualisations improve the cognitive understanding of complex climate tipping systems among non-specialists, compared with conventional static ones. It uses the Atlantic Meridional Overturning Circulation (AMOC) as a case study for a transferable framework built in three stages: \emph{detect}, a machine-learning analysis of the system's structure; \emph{translate}, an interactive artefact that renders that structure while holding its static and experiential conditions informationally equal; and \emph{evaluate}, an exploratory between-subjects study of twenty-five non-specialists, whose written responses were analysed against a pre-specified coding grid.

The results are informative precisely because they are not uniform. It has been found that experiential presentation did not always raise comprehension. It undoubtedly helped where a system's structure is temporal and invisible to a still figure, showed little to no advantage where a static figure already presented the structure plainly, and, for one scene, made a systematic misreading more likely rather than less. Participants themselves named the mechanism on both sides. Just to be clear, these results were tested against AMOC and its most commonly analysed and monitored plots, so the outcome is not meant to be transferred or taken as truth --- it is all about the framework.

Therefore, the contribution is mainly methodological: a framework, and a first body of evidence, for making invisible systems intelligible --- not by proving that experience helps everywhere, but by showing for which structures, for which viewers, and at what perceptual cost.

\cleardoublepage